\begin{abstractCN}
\setlength{\baselineskip}{23pt} %设置行距23磅

随着物联网的发展，网络中的状态更新业务呈现出终端数量多、短包突发等特点，此类业务已经广泛出现在工业物联网、车路协同等场景中。对这类系统而言，接收端更关心的是当前掌握的状态信息是否足够新，能否支撑后续的监测、估计与控制。传统的吞吐量和时延指标难以对信息的新鲜度进行衡量，为解决这个问题，信息年龄（Age of Information，AoI）被提出并受到广泛关注。

在海量终端的接入场景中，若持续依赖集中式调度，系统需要维护大量节点的状态，这将承担极高的控制和反馈开销，而且可扩展性往往也会受到限制。相比之下，随机接入则避免了对全局状态信息的掌控，能够在低复杂度的情况下支持大规模设备的动态接入，因此在拥有海量设备的通信场景中更加适用。正因如此，如何在大规模接入的条件下兼顾接入效率与信息新鲜度，是一个值得重点研究的 MAC 层问题。近年来，既有学者设计退避算法来规避碰撞，也有学者通过调节网络中节点的接入行为，以避免拥塞，从而来优化网络中的 AoI。总而言之，围绕 AoI 的随机接入机制研究已经取得了一定进展。

然而，现有机制没有充分考虑高并发条件下持续碰撞所带来的影响，而且大部分研究的讨论重点都放在平均 AoI 上，对 AoI 违规概率这类尾部指标的关注较少。对于工业控制等时间敏感型业务，仅考察平均指标不足以判断系统的可靠性。围绕这两个问题，本文设计了一种新的随机接入机制，并引入时隙 ALOHA 和 802.11 DCF 中进行性能分析，还专门对违规概率进行了研究。本文的主要工作如下：

\begin{enumerate} [wide]
    \item \textbf{设计面向大规模接入场景的等待机制。}针对 AoI 这一指标，本文提出了一种基于等待的随机接入机制，该机制将碰撞的客观事实考虑在内，补足了现有的优秀机制在碰撞处理方面的空缺。
    \item \textbf{等待 ALOHA 的理论分析。} 将等待机制引入时隙 ALOHA 框架，构建等待 ALOHA 协议，建立相应的理论模型，并推导平均 AoI 与 AoI 违规概率的闭式表达式，分析等待窗口和网络规模对系统性能的影响。
    \item \textbf{等待 CSMA 的理论分析。} 将等待机制进一步推广到 802.11 DCF 中，构建等待 CSMA 协议，并结合载波侦听、退避冻结和变长系统时隙等因素建立相应的理论模型，推导平均 AoI 与 AoI 违规概率的闭式表达式，分析等待窗口和网络规模对系统性能的影响。
    \item \textbf{揭示等待机制的作用规律。}基于闭式表达式并结合仿真验证，本文发现，在中高负载尤其是海量设备接入的场景下，适度等待并不会损害信息新鲜度。相反，它能够减少信道中的竞争，降低碰撞概率，从而同时改善平均 AoI 与 AoI 违规概率。
\end{enumerate}


\end{abstractCN}


\begin{keywordCN}
随机接入，信息年龄，性能分析
\end{keywordCN}

\begin{abstractEN}
\setlength{\baselineskip}{23pt} %设置行距23磅

With the development of the Internet of Things, especially massive machine-type communications (mMTC), status-update traffic in wireless networks has become characterized by a large number of terminals and bursty short-packet transmissions. Such traffic patterns are now common in scenarios such as industrial IoT and vehicle-road coordination. In these systems, the receiver is more concerned with whether the currently available state information is fresh enough to support subsequent monitoring, estimation, and control. Since traditional throughput and delay metrics cannot adequately capture this requirement, Age of Information (AoI) has been introduced to measure information freshness.

In massive-access scenarios, persistent reliance on centralized scheduling requires the system to maintain the states of a large number of nodes, which incurs very high control and feedback overhead and often limits scalability. By contrast, random access avoids maintaining global state information and can support dynamic access for a large number of devices with relatively low complexity, making it more suitable for communication scenarios with massive numbers of devices. For this reason, how to balance access efficiency and information freshness under large-scale contention has become an important MAC-layer problem. In recent years, some studies have designed backoff algorithms to mitigate collisions, while others have optimized AoI by regulating how nodes access the channel so as to avoid congestion. Overall, AoI-oriented random access mechanisms have made meaningful progress.

Nevertheless, existing mechanisms do not sufficiently account for the impact of persistent collisions under high concurrency, and most prior studies focus mainly on average AoI while paying much less attention to tail metrics such as the AoI violation probability. For time-sensitive applications such as industrial control, average performance alone is insufficient to assess system reliability. To address these two issues, this thesis designs a new random access mechanism, incorporates it into slotted ALOHA and IEEE 802.11 DCF, and analyzes its performance, with particular attention to the AoI violation probability. The main contributions are summarized as follows:

\begin{enumerate} [wide]
\item \textbf{Design of a wait-based mechanism for massive-access scenarios.} To optimize AoI, this thesis proposes a wait-based random access mechanism that explicitly takes collisions into account, thereby filling a gap in collision handling left by existing representative schemes.

\item \textbf{Theoretical analysis of Wait-ALOHA.} The proposed waiting mechanism is introduced into the slotted ALOHA framework to construct the Wait-ALOHA protocol. A corresponding theoretical model is established, and closed-form expressions for the average AoI and the AoI violation probability are derived. The effects of the waiting window and network size on system performance are also analyzed.

\item \textbf{Theoretical analysis of Wait-CSMA.} The proposed waiting mechanism is further extended to IEEE 802.11 DCF to construct the Wait-CSMA protocol. By incorporating carrier sensing, backoff freezing, and variable-length system slots, a corresponding theoretical model is established, and closed-form expressions for the average AoI and the AoI violation probability are derived. The effects of the waiting window and network size on system performance are also analyzed.

\item \textbf{Revealing the role of waiting.} Based on the closed-form expressions and supported by simulations, this thesis shows that under medium-to-high load, especially in massive-device access scenarios, moderate waiting does not degrade information freshness. On the contrary, it can reduce channel contention and lower the collision probability, thereby improving both the average AoI and the AoI violation probability.
\end{enumerate}

\end{abstractEN}

\begin{keywordEN}
Random access, Age of information, Performance analysis
\end{keywordEN}
